%% Formule de l'operateur : Stationnarite_Ustar
%% Dans ce fichier, les �ventuels accents sont cod�s
%% en ISO-8859-15 (aka latin9)
%% La ligne suivante donne trois exemples 
%% e accent aigu : �, e accent grave : �, e-dans-l'o : �

$M$ est le nombre de sous-p�riodes sur lesquelles on calcule $U_*$. 
Donc $M$ vaut $T_2 /T_1$ ( $T_2$ > $T_1$ et $T_2$ multiple de $T_1$ )\\

$U_*$ sur la $i$�me p�riode $T_1$ est calcul�e par :
$$
(U_*)_i = \left( \left ( \overline{U'W'} \right)^2_i + 
                 \left ( \overline{V'W'} \right)^2_i \right)^{1/4} 
$$

Puis $U_*$ moyen vaut
$$
U_{*moy} = \frac{1}{M} \sum_i (U_*)_i
$$
L'indice de stationnarit� vaut :
$$
Istat = \left\vert  \frac{ U_*  - U_{*moy} }
                   { U_* } \right\vert
$$
avec $U_*$ : vitesse de frottement calcul�e sur toute la p�riode $T_2$